% from ega2-1-fr.tex
\oldpage[II]{80}

On sait (\hyperref[II.3.4.7-fr]{3.4.7}) que l'existence d'un homomorphisme
surjectif $\varphi:q^*(\mathcal{E})\to\mathcal{L}$ ($\mathcal{E}$ étant un
$\mathcal{O}_Y$-Module quasi-cohérent) entraîne que $\sigma$ est surjectif~;
en outre, à la factorisation
$q^*(\mathcal{E})\to q^*(q_*(\mathcal{L}))\xrightarrow{\sigma}\mathcal{L}$
de $\varphi$ correspond canoniquement une factorisation
\[
  q^*(\mathbf{S}(\mathcal{E}))
    \longrightarrow q^*(\mathbf{S}(q_*(\mathcal{L})))
    \longrightarrow\bigoplus_{n\geq0}\mathcal{L}^{\otimes n}
\]
donc (\hyperref[II.3.8.3-fr]{3.8.3}) l'hypothèse que
$r_{\mathcal{L},\varphi}$ est une immersion entraîne qu'il en est de même de
$j=r_{\mathcal{L},\sigma}$~; en outre
(\hyperref[II.4.2.4-fr]{4.2.4}), $\mathcal{L}$ est isomorphe à
$j^*(\mathcal{O}_{P'}(1))$ en posant
$P'=\mathbf{P}(q_*(\mathcal{L}))$. On voit donc que a) et b) sont
équivalents.
\end{proof}

\begin{corollary}[4.4.5]
\label{II.4.4.5-fr}
Supposons $q$ quasi-compact. Pour que $\mathcal{L}$ soit très ample
relativement à $Y$, il faut et il suffit qu'il existe un recouvrement ouvert
$(U_\alpha)$ de $Y$ tel que
$\mathcal{L}|q^{-1}(U_\alpha)$ soit très ample relativement à $U_\alpha$ pour
tout $\alpha$.
\end{corollary}

\begin{proof}
En effet, la condition b) de (\hyperref[II.4.4.4-fr]{4.4.4}) est locale sur
$Y$.
\end{proof}

\begin{proposition}[4.4.6]
\label{II.4.4.6-fr}
Soient $Y$ un schéma quasi-compact, ou un préschéma dont l'espace sous-jacent
est noethérien, $q:X\to Y$ un morphisme de type fini, $\mathcal{L}$ un
$\mathcal{O}_X$-Module inversible. Les conditions a) et b) de
(\hyperref[II.4.4.4-fr]{4.4.4}) sont alors équivalentes aux suivantes~:
\begin{enumerate}
  \item[\rm{a')}] Il existe un $\mathcal{O}_Y$-Module quasi-cohérent de type
  fini $\mathcal{E}$ et un homomorphisme surjectif
  $\varphi:q^*(\mathcal{E})\to\mathcal{L}$ tel que
  $r_{\mathcal{L},\varphi}$ soit une immersion.
  \item[\rm{b')}] Il existe un sous-$\mathcal{O}_Y$-Module cohérent de type
  fini $\mathcal{E}$ de $q_*(\mathcal{L})$ ayant les propriétés énoncées dans
  a').
\end{enumerate}
\end{proposition}

\begin{proof}
Il est clair que a') ou b') implique a)~; d'autre part a) implique a') en
vertu de (\hyperref[II.4.4.1-fr]{4.4.1}) et de même b) implique b').
\end{proof}

\begin{corollary}[4.4.7]
\label{II.4.4.7-fr}
Supposons que $Y$ soit un schéma quasi-compact ou un préschéma noethérien. Si
$\mathcal{L}$ est très ample pour $q$, il existe une $\mathcal{O}_Y$-Algèbre
graduée quasi-cohérente $\mathcal{S}$, telle que $\mathcal{S}_1$ soit de type
fini et engendre $\mathcal{S}$, et une $Y$-immersion ouverte dominante
$i:X\to P=\operatorname{Proj}(\mathcal{S})$ telle que $\mathcal{L}$ soit
isomorphe à $i^*(\mathcal{O}_P(1))$.
\end{corollary}

\begin{proof}
En effet, la condition b') de (\hyperref[II.4.4.6-fr]{4.4.6}) est vérifiée~;
le morphisme structural $p:\mathbf{P}(\mathcal{E})=P'\to Y$ est alors séparé
et de type fini (\hyperref[II.3.1.3-fr]{3.1.3}), donc $P'$ est un schéma
quasi-compact (resp. un préschéma noethérien) si $Y$ est un schéma
quasi-compact (resp. un préschéma noethérien). Soit $Z$ l'adhérence
(\hyperref[I.9.5.11-fr]{I, 9.5.11}) du sous-préschéma $X'$ de $P'$ associé à
l'immersion $j=r_{\mathcal{L},\varphi}$ de $X$ dans $P'$~; il est clair que
$j$ se factorise en l'injection canonique $Z\to P'$ et en une immersion
ouverte dominante $i:X\to Z$. Mais $Z$ s'identifie à un préschéma
$\operatorname{Proj}(\mathcal{S})$, où $\mathcal{S}$ est une
$\mathcal{O}_Y$-Algèbre graduée quotient de $\mathbf{S}(\mathcal{E})$ par un
faisceau quasi-cohérent gradué d'idéaux
(\hyperref[II.3.6.2-fr]{3.6.2}), et il est clair que $\mathcal{S}_1$ est de
type fini et engendre $\mathcal{S}$~; en outre $\mathcal{O}_Z(1)$ est l'image
réciproque de $\mathcal{O}_{P'}(1)$ par l'injection canonique
(\hyperref[II.3.6.3-fr]{3.6.3}), donc
$\mathcal{L}=i^*(\mathcal{O}_Z(1))$.
\end{proof}

\begin{proposition}[4.4.8]
\label{II.4.4.8-fr}
Soient $q:X\to Y$ un morphisme, $\mathcal{L}$ un $\mathcal{O}_X$-Module très
ample relativement à $q$, $\mathcal{L}'$ un $\mathcal{O}_X$-Module inversible
tel qu'il existe un $\mathcal{O}_Y$-Module quasi-cohérent $\mathcal{E}'$ et un
homomorphisme surjectif $q^*(\mathcal{E}')\to\mathcal{L}'$. Alors
$\mathcal{L}\otimes_{\mathcal{O}_X}\mathcal{L}'$ est très ample relativement
à $q$.
\end{proposition}

\begin{proof}
En effet, l'hypothèse entraîne l'existence d'un $Y$-morphisme
$r':X\to P'=\mathbf{P}(\mathcal{E}')$ tel que
$\mathcal{L}'=r'^*(\mathcal{O}_{P'}(1))$
(\hyperref[II.4.2.1-fr]{4.2.1}). Il y a par hypothèse un
$\mathcal{O}_Y$-Module quasi-cohérent $\mathcal{E}$ et une

